% OLD TEXT FROM HERE ON
% This chapter contains text from the original submission.
% It will be revised as part of the salvage plan.

\section{Bridging paradigms in CGE modeling? A preliminary assessment} \label{sec:conclusion}

In sum, this paper observes and analyzes two axiomatic variations of the standard neoclassical approach towards input-output modeling. The first relates to the more long-standing tradition of modeling some substitution processes as rather inflexible, which can be mimicked by the use of a CES-framework with appropriate parameter values. In this case the underlying CES function allows for establishing a plausible theoretical continuum between different paradigmatic perspectives on the level of assumption, but when applied in isolation does not contribute to closing paradigmatic gaps in terms of outcomes. Moreover, as indicated by Figure \ref{fig:labour_camp}, which summarizes outcomes from all scenarios analyzed, variations in the imposed elasticity have little to no impact on final results for large parts of the relevant parameter space.

\begin{figure}
    \centering
    \includegraphics[width=1\linewidth]{pictures/comparison_between_labor_slacks.png}
    \caption{Different estimates of the impact of shock on final demand focused on investments towards a carbon-neutral residential sector. Note that, in the visual display there is a strong overlay between all scenarios associated with using a CES-function for the case of calibrated labor slack as already shown in Figure \ref{fig:elasticity-gradient2}.}
    \label{fig:labour_camp}
    \end{figure}

The second variation -- the introduction of labor slack -- on the other hand is useful for creating outcomes that align more closely with those achieved by the traditional Leontief approach. However, the assumption of positive labor slack in this context plays the role of a \textit{deus ex machina} that implicitly governs the GDP-expansion achieved as increases in input map roughly one-to-one on increases in aggregate output. Hence, on this level a fundamental cleavage in terms of assumptions remains as the Leontief model conceives new equilibria as essentially demand-driven, which is a feature that cannot be reproduced in a general equilibrium context by simply adding slack on factor markets. Rather, in this framework equilibria are driven by supply-considerations, which are modified ad-hoc to provide alternative outcomes. This latter aspect as is also evident from the fact that typical output expansions predicted in Figure \label{fig:labour_camp} for the relevant scenarios perfectly match the respective amount of labor slack, i.e., the increase in labor inputs. This latter feature makes it hard to achieve a common ground as the effects of a demand stimulus on exact allocation of labor remains fundamentally contested.

In terms of policy-relevance our results indicate that especially the latter variation of introducing labor slack, opens some leeway to generate intermediate results, that can be situated in-between traditional approaches. However, for doing so one must somehow specify, how many hitherto unutilized factor resources can be mobilized in case of an exogenous demand shock, as well as accept that these assumptions on additional factor resources -- their total amount as well as their allocation across sector --, will play a more decisive role in coining final results, than the exogenous demand shock, that is employed as an original stimulus.

